% =======================================================
% 1. INTRODUZIONE
% =======================================================
\section{Introduzione}
Il presente documento costituisce il Piano di Progetto (PdP) relativo allo sviluppo del sistema Second Brain. Esso definisce l'organizzazione del lavoro, i rischi identificati, le risorse disponibili, la suddivisione delle attività e il calendario delle stesse, in conformità con le pratiche di ingegneria del software.

\subsection{Scopo del documento}
Il Piano di Progetto relativo allo sviluppo del capitolato Second Brain ha i seguenti obbiettivi:
\begin{itemize}
    \item definire la struttura organizzativa e i ruoli del team
    \item identificare e gestire i rischi in modo proattivo
    \item fornire una visione chiara e aggiornata dell'andamento del progetto
    \item pianificare le attività nel breve e nel lungo termine
    \item la stima dei costi e il consuntivo periodico
\end{itemize}

\subsection{Scopo del prodotto}
Zucchetti propone la realizzazione di una piattaforma web orientata alla scrittura e alla gestione avanzata di contenuti testuali, concepita come un ambiente intelligente di supporto all’utente. Il sistema si baserà su Markdown per la formattazione dei documenti e integrerà modelli linguistici di grandi dimensioni (LLM) per ampliare le funzionalità disponibili.
All’interno dell’applicazione, l’utente potrà redigere testi con aggiornamento immediato della visualizzazione, oltre a sfruttare strumenti automatizzati per la produzione e l’elaborazione dei contenuti, come generazione, sintesi, riformulazione e traduzione. Sarà inoltre possibile adottare approcci strutturati all’analisi, tra cui la tecnica dei “Sei cappelli per pensare”, utile per valutare un testo da molteplici prospettive. Un’ulteriore caratteristica consisterà nell’utilizzo di istruzioni testuali per guidare la creazione di nuovi contenuti, secondo il paradigma del cosiddetto Distant Writing.
Nel complesso, la soluzione mira a fornire un ambiente evoluto per la scrittura creativa e l’organizzazione delle idee, sfruttando le capacità degli LLM per migliorare il processo di produzione e revisione dei testi.

\subsection{Glossario}
\url{{ https://synergo-unit-unipd.github.io/Documentary-Repo/docs/RTB/doc_gestionali/doc_interni/glossario/glossario.pdf}}


\subsection{Riferimenti}
\subsubsection{Normativi}
\begin{itemize}
	\item Capitolato d'appalto: \url{{https://www.math.unipd.it/~tullio/IS-1/2025/Progetto/C6.pdf}}
	\item \textit{Norme di Progetto v1.1.2}
	\item Regolamento del Progetto Didattico:\url{{ https://www.math.unipd.it/~tullio/IS-1/2009/Approfondimenti/ISO_12207-1995.pdf}}
\end{itemize}

\subsubsection{Informativi}
\begin{itemize}
	\item \textit{Analisi dei Requisiti vX.X.X}
	\item Verbali degli incontri interni ed esterni
	\item Slide del corso di Ingegneria del Software (T04 - Gestione di progetto)
\end{itemize}

\subsection{Scadenze}
Le consegne concordate per le revisioni di progetto sono le seguenti:
\begin{itemize}
	\item \textbf{RTB (Requirements and Technology Baseline):} {Marzo-Luglio}
	\item \textbf{PB (Product Baseline):} {Luglio-Agosto}
\end{itemize}